\documentclass[10pt,letterpaper,oneside]{report}
\usepackage[T1]{fontenc}
\usepackage[utf8]{inputenc}
\usepackage{lmodern}
\usepackage{microtype}
\usepackage[margin=0.85in,headheight=15pt]{geometry}
\usepackage{enumitem,booktabs,longtable,tabularx,array,makecell}
\usepackage{fancyhdr,lastpage,xcolor,hyperref,titlesec}
\usepackage{listings}
\usepackage[most]{tcolorbox}
\definecolor{ink}{HTML}{182D40}
\definecolor{accent}{HTML}{126B77}
\definecolor{pale}{HTML}{EEF5F5}
\definecolor{grayblue}{HTML}{677887}
\definecolor{codebg}{HTML}{F5F7F8}
\hypersetup{colorlinks=true,linkcolor=accent,urlcolor=accent,pdfauthor={C23 Engineering Companion},pdfsubject={ISO/IEC 9899:2024 educational engineering companion}}
\titleformat{\chapter}[display]{\normalfont\huge\bfseries\color{ink}}{\filleft\large\color{accent}CHAPTER \thechapter}{1ex}{\titlerule[1pt]\vspace{1ex}}
\titleformat{\section}{\Large\bfseries\color{ink}}{\thesection}{0.7em}{}
\titleformat{\subsection}{\large\bfseries\color{accent}}{\thesubsection}{0.65em}{}
\setlist[itemize]{leftmargin=1.3em,itemsep=2pt,topsep=3pt}
\setlist[enumerate]{leftmargin=1.6em,itemsep=3pt,topsep=3pt}
\setlength{\parskip}{5pt}
\setlength{\parindent}{0pt}
\renewcommand{\arraystretch}{1.18}
\pagestyle{fancy}\fancyhf{}\fancyhead[L]{\small\sffamily C23 ENGINEERING SERIES}\fancyhead[R]{\small\sffamily\color{accent}\DocShort}\fancyfoot[L]{\small\color{grayblue}ISO/IEC 9899:2024 companion}\fancyfoot[R]{\small\thepage\ / \pageref{LastPage}}
\newtcolorbox{keybox}{breakable,colback=pale,colframe=accent,boxrule=.6pt,arc=1mm,left=8pt,right=8pt,top=6pt,bottom=6pt}
\newtcolorbox{warnbox}{breakable,colback=codebg,colframe=grayblue,boxrule=.6pt,arc=1mm,left=8pt,right=8pt,top=6pt,bottom=6pt}
\lstdefinestyle{cstyle}{language=C,basicstyle=\ttfamily\footnotesize,keywordstyle=\color{accent}\bfseries,commentstyle=\color{grayblue}\itshape,stringstyle=\color{ink},backgroundcolor=\color{codebg},frame=single,rulecolor=\color{grayblue!40},breaklines=true,showstringspaces=false,tabsize=4,columns=fullflexible,keepspaces=true,upquote=true}
\lstset{style=cstyle}
\newcommand{\refcl}[1]{\textsf{\footnotesize[ISO/IEC 9899:2024, #1]}}
\newcommand{\standard}{\textbf{[STANDARD]}}
\newcommand{\optional}{\textbf{[OPTIONAL]}}
\newcommand{\guidance}{\textbf{[GUIDANCE]}}
\newcommand{\testtag}{\textbf{[TEST]}}
\newcommand{\term}[1]{\textbf{#1}}
\newcommand{\DocShort}{LANGUAGE REFERENCE}
\begin{document}
\begin{titlepage}
\color{ink}\vspace*{1.2in}
{\sffamily\small\color{accent}\bfseries TECHNICAL DOCUMENTATION SUITE \hfill EDITION 2026-10}\par
\vspace{.3in}\rule{\linewidth}{1.5pt}\par\vspace{.22in}
{\Huge\bfseries C23 Language Technical Reference}\par\vspace{.22in}
{\Large\color{grayblue}Syntax, semantics, objects, expressions, declarations and preprocessing}\par\vspace{.4in}
\begin{keybox}\textbf{Standards basis:} ISO/IEC 9899:2024 (C23). Independent educational and engineering companion. Not an official ISO interpretation, full specification, certification or conformance assessment.\end{keybox}
\vfill{\large\bfseries C23 Software Engineering and Language Standards}\par\vspace{.15in}{\color{grayblue}Prepared October 2026 \quad | \quad Source edition: October 2024}\par
\end{titlepage}
\pagenumbering{roman}\tableofcontents\clearpage\pagenumbering{arabic}

\chapter*{Sources and interpretation key}
\addcontentsline{toc}{chapter}{Sources and interpretation key}
\textbf{Primary source:} ISO/IEC 9899:2024(en), \emph{Information technology --- Programming languages --- C}, fifth edition (2024). Clause and annex identifiers refer to the supplied licensed standard. This handbook contains original explanations, designs and examples; it does not reproduce the standard. The licensed standard remains the authoritative source for exact wording and conformance decisions.\par
\textbf{Secondary engineering sources:} \url{https://gcc.gnu.org/onlinedocs/gcc/Standards.html}; \url{https://clang.llvm.org/c_status.html}; \url{https://cmake.org/cmake/help/latest/prop_tgt/C_STANDARD.html}. Compiler support and diagnostics are release-specific: verify the installed compiler, target and library before relying on a feature.\par
\begin{keybox}\textbf{Interpretation categories.} \textbf{[STANDARD]} summarizes a requirement, constraint, semantic rule or defined facility and cites the applicable clause. \textbf{[OPTIONAL]} describes conditional or implementation-dependent support. \textbf{[GUIDANCE]} records an engineering recommendation, not an ISO mandate. \textbf{[TEST]} is an original check or exercise; a passing test is evidence but is not proof of full conformity.\end{keybox}
\textbf{Copyright boundary:} The supplied PDF is individually licensed. Do not redistribute it with this suite; distribute only these independently written companion materials in accordance with applicable rights.\par


\chapter{Conformance and the interpretation of C}
\section{A programming-language standard, not a secure-coding standard}
\standard\ The scope covers the representation and interpretation of C programs, program input and output, and implementation restrictions. It does not prescribe a build system, operating-system API, secure-coding process or a particular compiler architecture. \refcl{1}. A conforming implementation and a strictly conforming program are not interchangeable terms. \refcl{4}.
\section{How to read clauses}
The document distinguishes \emph{Syntax}, \emph{Constraints}, \emph{Semantics}, \emph{Description}, \emph{Returns}, \emph{Runtime-constraints}, \emph{Environmental limits}, and \emph{Recommended practice}. A syntax or constraint violation triggers translation diagnostics under the specified conditions; undefined behavior elsewhere need not produce one. \refcl{Introduction, 5.2.1.3}.
\begin{keybox}\standard\ An implementation diagnoses at least one syntax or constraint violation in a preprocessing translation unit or translation unit; successful compilation alone does not establish a program's validity. \refcl{5.2.1.3}.\end{keybox}
\section{Behavioral vocabulary}
\begin{longtable}{p{.28\linewidth}p{.65\linewidth}}\toprule Term & Consequence for design and review\\\midrule\endhead
Undefined behavior & The standard imposes no requirements in the triggering situation; never treat a specific observed result as a contract. \refcl{3.5.3}\\
Unspecified behavior & One of the permitted alternatives may be selected without an obligation to document the choice. \refcl{3.5.4}\\
Implementation-defined behavior & A choice exists, and the implementation documents how it is made. \refcl{3.5.1}\\
Locale-specific behavior & Depends on documented locale conventions. \refcl{3.5.2}\\
Bounded / critical undefined behavior & C23 introduces a distinction based on potential out-of-bounds stores; it does not make either kind safe. \refcl{3.30--3.31; Annex L}\\
\bottomrule\end{longtable}
\chapter{Environment and translation model}
\section{Translation phases}
\standard\ Source character mapping, line splicing, tokenization, preprocessing, character conversion, literal concatenation, syntactic analysis and linkage are conceptual stages. An implementation need not literally execute eight separate passes. \refcl{5.2.1.1}.
\guidance\ Run both a preprocessing-only inspection (e.g., \texttt{-E}) and an actual translation; preprocessor success does not demonstrate semantic correctness.
\section{Hosted and freestanding}
\standard\ The standard recognizes hosted and freestanding execution environments, with different startup and library obligations. \refcl{5.2.2}.\guidance\ Declare the intended environment in the toolchain manifest: target triple, ABI, runtime, libc, linker, compiler and flags. Avoid assuming a hosted header is available on a bare-metal target.
\section{Execution model}
Side effects, evaluations, object lifetimes and signal or thread interactions make some source-level transformations invalid even when arithmetic identities appear sound. \refcl{5.1.2, 5.2.3, 5.3.4}.
\chapter{Types, objects, storage, and representations}
\section{Scope, linkage, lifetime and storage duration}
Separate lexical visibility (scope), entity identity across declarations (linkage), storage duration, and the valid lifetime of an object. \refcl{6.2.1--6.2.4}. Automatic and allocated storage differ; a pointer may remain in scope after the pointed-to object's lifetime has ended.
\begin{lstlisting}
int *bad_pointer(void) {
    int local = 7;
    return &local; /* invalid after local's lifetime ends */
}
\end{lstlisting}
\guidance\ Every pointer-bearing API should document ownership, lifetime, permitted aliasing, mutability and alignment.
\section{Type families and composite types}
Arithmetic, integer, floating, pointer, array, structure, union, function, atomic and qualified types carry distinct rules. Declarations of an entity across translation units must satisfy compatibility and linkage requirements. \refcl{6.2.5--6.2.7}.
\section{Object representations and alignment}
Reading an object's byte representation is distinct from interpreting arbitrary bytes as a valid typed value. Representation, effective types, alignment and the applicable aliasing rules constrain optimization and portability. \refcl{6.2.6, 6.2.8}.
\begin{lstlisting}
#include <string.h>
#include <stdint.h>
uint32_t read_word(const unsigned char bytes[4]) {
    uint32_t value = 0;
    memcpy(&value, bytes, sizeof value); /* host representation */
    return value;
}
\end{lstlisting}
\guidance\ For interchange formats, assemble integers from explicitly ordered octets rather than relying on host endianness or object representation.
\chapter{Arithmetic, conversions, and expressions}
\section{Integer promotions and usual arithmetic conversions}
\standard\ Integer promotions and the usual arithmetic conversions determine operand types and the result type; explicit casts can change the result without preventing overflow. \refcl{6.3.1, 6.5}.
\begin{lstlisting}
#include <limits.h>
#include <stdbool.h>
bool checked_add_int(int a, int b, int *out) {
    if ((b > 0 && a > INT_MAX - b) ||
        (b < 0 && a < INT_MIN - b)) return false;
    *out = a + b;
    return true;
}
\end{lstlisting}
\section{Sequencing and side effects}
\standard\ A full expression introduces sequencing boundaries; not all operands are evaluated in textual order. Mixing modifications and accesses of the same scalar object without adequate sequencing can cause undefined behavior. \refcl{6.5.1, 6.8}.
\begin{lstlisting}
int i = 1;
/* Avoid constructs that depend on unsequenced side effects. */
int a = i;
++i;
int b = i;
\end{lstlisting}
\section{Pointers and arrays}
An array expression often converts to a pointer to its first element, but pointer arithmetic remains constrained to the array object and the one-past position. One-past pointers cannot be dereferenced. \refcl{6.3.3, 6.5.7}.
\section{Selected operators}
Unary, multiplicative, additive, shift, relational, equality, bitwise, logical, conditional, assignment and comma operators have distinct constraints and sequencing. \refcl{6.5.2--6.5.18}.\guidance\ Prefer explicit parentheses where a reviewer would otherwise need a precedence table.
\chapter{Lexical structure, identifiers and constants}
\section{Tokens, source encoding and names}
\standard\ Keywords, identifiers, constants, string literals, punctuators and header names are specified in \refcl{6.4}; Unicode-based identifier rules and character mappings have portability implications. Compile the source with a declared input encoding rather than trusting platform defaults.
\section{Literals}
C23 provides binary integer constants and digit separators. The value, not the punctuation, participates in arithmetic. \refcl{6.4.5; Annex M.2}.
\begin{lstlisting}
unsigned permissions = 0b1101'0010;
/* Same value as 0xD2, independent of display base. */
\end{lstlisting}
\section{Strings and encodings}
Distinguish a \texttt{char} array from a null-terminated byte string. Literal encodings, Unicode encodings, and runtime locale conversions must not be conflated. \refcl{6.4.6, 6.2.9, 7.30}.
\chapter{Declarations and the C23 type system}
\section{Declarators and qualifiers}
\texttt{const}, \texttt{volatile}, \texttt{restrict}, and atomic qualification are not interchangeable. Use \texttt{restrict} only with a valid nonaliasing contract; \texttt{volatile} is not a threading synchronization primitive. \refcl{6.7.4, 6.7.7, 7.17}.
\section{Type inference, \texttt{constexpr} and type queries}
\standard\ C23 expands \texttt{auto} inference, adds \texttt{typeof} and \texttt{typeof\_unqual}, and introduces \texttt{constexpr} for qualifying object definitions under specified constraints. \refcl{6.7.2, 6.7.10, Annex M.2}.
\begin{lstlisting}
/* Compile only after checking these C23 features in the toolchain. */
constexpr int buffer_capacity = 64;
auto counter = 0;
typeof(counter) another = 1;
static_assert(buffer_capacity >= 32);
\end{lstlisting}
\guidance\ Explicit public API types often communicate ABI intent better than inferred types; favor inference for well-scoped locals.
\section{Bit-precise integers and enumerations}
\standard\ The \texttt{\_BitInt(N)} family represents implementation-supported bit-precise integer widths subject to the standard's limits; C23 also adds fixed underlying enumeration types. \refcl{6.7.3, Annex M.2}.\optional\ Availability of a requested width and ABI behavior require compiler and target checks.
\begin{lstlisting}
/* Feature-probe before enabling in production. */
unsigned _BitInt(17) field = 0;
enum packet_kind : unsigned char { PKT_DATA = 1, PKT_ACK = 2 };
\end{lstlisting}
\section{Initialization, empty braces and attributes}
Aggregate initialization, designated initialization and C23 empty initialization should be evaluated in context. Attributes annotate constructs; the standard makes support for standard attributes optional and implementation-defined, so use \texttt{\_\_has\_c\_attribute} when feature-dependent. \refcl{6.7.11--6.7.13}.
\begin{lstlisting}
struct Sample { int x; int y; };
struct Sample s = {};  /* C23 empty initialization */
#if defined(__has_c_attribute)
# if __has_c_attribute(nodiscard)
[[nodiscard]]
# endif
#endif
int parse_result(const char *text);
\end{lstlisting}
\chapter{Statements, functions and program structure}
\section{Control flow}
Blocks, selection, iteration, jumps and labels are specified in \refcl{6.8}. C23 permits labels before declarations and at the end of blocks; that is not a recommendation to use \texttt{goto} for ordinary control flow.
\section{Function interfaces}
\standard\ An empty parameter list declares a function with no parameters in C23; identifier-list (K\&R) definitions are no longer supported. \refcl{6.7.7, 6.9.1, Annex M.2}.
\begin{lstlisting}
int next_sequence(void); /* explicit no-parameter contract */
int next_sequence() { return 1; } /* same contract in C23 */
\end{lstlisting}
\guidance\ Header declarations should consistently use \texttt{(void)} where cross-edition compatibility matters.
\chapter{The preprocessor}
\section{Conditional compilation and replacement}
Includes, macros, conditional inclusion, diagnostic directives, predefined macros and \texttt{\_Pragma} have different processing rules. \refcl{6.10}. Use include guards even where a compiler has a once-only optimization.
\begin{lstlisting}
#ifndef EXAMPLE_H
#define EXAMPLE_H
#if defined(__has_include)
# if __has_include(<stdckdint.h>)
#  include <stdckdint.h>
# endif
#endif
#endif
\end{lstlisting}
\section{C23 directives and expansion}
\standard\ C23 adds \texttt{\#elifdef}, \texttt{\#elifndef}, \texttt{\#warning}, \texttt{\#embed}, and enhanced variadic macro control with \texttt{\_\_VA\_OPT\_\_}. \refcl{6.10.2--6.10.7; Annex M.2}.
\begin{lstlisting}
#if defined(PLATFORM_A)
# define PLATFORM_NAME "A"
#elifdef PLATFORM_B
# define PLATFORM_NAME "B"
#else
# warning "Unknown platform"
#endif
\end{lstlisting}
\optional\ \texttt{\#embed} and associated parameters must be tested on the actual compiler; the directive's presence does not make path resolution or build artifacts portable by itself.
\chapter{Concurrency and the C memory model}
\section{Atomic data and threads}
Use \texttt{<stdatomic.h>} operations to express ordering and visibility; use \texttt{<threads.h>} for the standardized thread interfaces where the implementation supplies them. A non-atomic data race has undefined behavior. \refcl{5.1.2.4, 7.17, 7.28}.
\begin{lstlisting}
#include <stdatomic.h>
atomic_bool ready = false;
/* Producer: publish after preparing shared state. */
void publish(void) { atomic_store_explicit(&ready, true,
                                          memory_order_release); }
/* Consumer: acquire before examining that state. */
int published(void) { return atomic_load_explicit(&ready,
                                        memory_order_acquire); }
\end{lstlisting}
\guidance\ For compound invariants across several variables, a mutex is usually easier to prove correct than a hand-built lock-free protocol.
\chapter{Language portability checklist}
\begin{enumerate}
\item Document data-model assumptions: integer widths, byte size, \texttt{char} signedness, alignment and endian representation. \refcl{5.3.5, 6.2.6; Annex J}.
\item Record expected evaluation/overflow behavior and eliminate undefined behavior, not just warning messages. \refcl{6.3, 6.5}.
\item Verify conformance mode, hosted/freestanding runtime, attributes and optional features on every target. \refcl{4, 5.2.2, 6.7.13}.
\item Keep library contracts, ownership and locale assumptions visible at every interface. \refcl{7.1, 7.11}.
\item Save a compiler-specific implementation-defined behavior register and negative tests. \refcl{Annex J.3}.
\end{enumerate}

\end{document}
